% 原书分章PDF第13—15页：第29.5节及完整表29.1。
\section{模拟}

我们考虑数据结构 $O=(W,A,Y)$ 的三个数据生成分布 $P_{0,1}$、$P_{0,2}$ 和 $P_{0,3}$。这三个分布，仅在给定 $(A,W)$ 时 $Y$ 的条件均值和方差方面存在差异。具体而言，从 $P_{0,j}$（$j=1,2,3$）中抽取的 $O=(W,A,Y)$ 满足：
\begin{itemize}
\item $W\equiv(V,W_1,W_2)$，其中 $P_{0,j}(V=1)=1/6$、$P_{0,j}(V=2)=1/3$、$P_{0,j}(V=3)=1/2$；给定 $V$，$(W_1,W_2)$ 是一个高斯随机向量。当 $V=1$ 时，均值为 $(0,0)$，方差为 $\left(\begin{smallmatrix}1&-0.2\\-0.2&1\end{smallmatrix}\right)$；当 $V=2$ 时，均值为 $(1,1/2)$，方差为 $\left(\begin{smallmatrix}0.5&0.1\\0.1&0.5\end{smallmatrix}\right)$；当 $V=3$ 时，均值为 $(1/2,1)$，方差为 $\left(\begin{smallmatrix}1&0\\0&1\end{smallmatrix}\right)$。
\item 给定 $W$，若 $W_1\ge1.1$ 且 $W_2\ge0.8$，则 $A=0$ 的概率为80\%，否则为10\%。此外，给定 $W$ 且 $A\ne0$，$3A-1$ 来自自由度为1、非中心参数为 $\sqrt{(W_1-1.1)^2+(W_2-0.8)^2}$ 的 $\chi^2$ 分布。
\item 给定 $(W,A)$，$Y$ 是一个高斯随机变量，其均值为
\begin{itemize}[label=--]
\item 对 $j=1,2$，为 $A(W_1+W_2)/6+W_1+W_2/4+\exp((W_1+W_2)/10)$；
\item 对 $j=3$，为 $A(W_1+W_2)/6+W_1+W_2/4+\exp((W_1+W_2)/10)+3AV$；
\end{itemize}
标准差为
\begin{itemize}[label=--]
\item 对 $j=1$，当 $V=1$、$V=2$、$V=3$ 时，分别为2、1.5和1；
\item 对 $j=2,3$，当 $V=1$、$V=2$、$V=3$ 时，分别为9、4和1。
\end{itemize}
\end{itemize}
真实参数近似等于 $\Psi^c(P_{0,1})=\Psi^c(P_{0,2})=0.1201$，以及 $\Psi^c(P_{0,3})=6.9456$。

取 $B=10^3$，对每个 $j=1,2,3$，独立重复以下步骤：
\begin{enumerate}
\item 模拟一个由从 $P_{0,j}$ 中独立抽取的 $N=10^7$ 个观测组成的数据集；
\item 通过简单随机抽样（simple random sampling，SRS；这与使用 $h_0\equiv1$ 的 $P^{s\mathrm P}$ 相同），从数据集中抽取 $n_0\equiv10^3$ 个观测，并根据这些观测：
\begin{enumerate}[label=(\alph*)]
\item 应用第29.4节描述的程序，取得 $f_{n_0,1}\equiv D^c(P_{n_0}^{k_{n_0}})$；
\item 将 $f_{n_0,1}(O)$ 和 $f_{n_0,1}(O)^2$ 对 $V$ 作回归，将由此得到的 $f_2^{\mathrm P}$ 估计的平方根称为 $f_{n_0,2}^{\mathrm P}$，将由此得到的 $f_2^\Pi$ 估计的平方根称为 $f_{n_0,2}^{\Pi}$，见式 (29.11)；
\item 估计 $V$ 的边际分布，用 $\pi_{n_0,2}$ 估计 $P_0f_{n_0,2}^{\mathrm P}$，并设定 $h_{n_0}^{\mathrm P}\equiv f_{n_0,2}^{\mathrm P}/\pi_{n_0,2}$、$h_{n_0}^{\Pi}\equiv f_{n_0,2}^{\Pi}/P_Nf_{n_0,2}^{\Pi}$，见式 (29.12) 和 (29.13)；
\end{enumerate}
\item 依次对每个 $n\in\{10^3,5\times10^3,10^4\}$，以及SRS、使用 $h_{n_0}^{\mathrm P}$ 的 $P^{s\mathrm P}$、使用 $P_{n_0}^{\Pi}$ 的 $P^{s\Pi}$ 这三种调查抽样设计中的每一种，从已经去除步骤2所抽观测的数据集中，抽取一个含有 $n$ 个观测的子样本，并根据这些观测应用第29.4节描述的程序。我们用式 (29.4) 给出的 $\hat\sigma_n^2$ 估计 $\sigma_1^2$，尽管事先不能确定它是一个保守估计量。
\end{enumerate}
由此得到 $27\times B$ 个估计及各自的CI。

为说明各情形中的 $h_{n_0}^{\mathrm P}$ 和 $h_{n_0}^{\Pi}$，我们报告它们在 $P_{0,1}$、$P_{0,2}$ 和 $P_{0,3}$ 下，跨 $B$ 次模拟研究的平均值：
\begin{itemize}[label=--]
\item 在 $P_{0,1}$ 下，预期 $h_{n_0}^{\mathrm P}$ 和 $h_{n_0}^{\Pi}$ 相似，实际也确实得到，它们平均而言近似等于 $(h_1(1),h_1(2),h_1(3))\approx(2.10,0.83,0.75)$；
\item 在 $P_{0,2}$ 下，同样预期 $h_{n_0}^{\mathrm P}$ 和 $h_{n_0}^{\Pi}$ 相似，实际也确实得到，它们平均而言近似等于 $(h_1(1),h_1(2),h_1(3))\approx(3.39,0.83,0.32)$；
\item 在 $P_{0,3}$ 下，并不预期 $h_{n_0}^{\mathrm P}$ 和 $h_{n_0}^{\Pi}$ 相似，所得平均值分别近似等于 $(h_1(1),h_1(2),h_1(3))\approx(2.93,0.66,0.58)$ 和 $(h_1(1),h_1(2),h_1(3))\approx(2.97,0.68,0.56)$。尽管差异很小，但具有显著性。
\end{itemize}

在R包 \texttt{tmle.npvi} 中，应用这一TMLE程序很直接（Chambaz 和 Neuvial，2016，2015）。不过，应注意，需要计算 $\hat\sigma_n^2$。具体而言，我们将 \texttt{iter}，即靶向步骤的最大迭代次数，设为 \texttt{7}，并将 \texttt{stoppingCriteria} 设为 \texttt{list(mic=0.01, div=0.001, psi=0.05)}，以此精细调整TMLE程序。此外，使用名为 \texttt{"learning"} 的默认 \texttt{flavor}，因而特别依赖参数线性模型，来估计无限维参数 $Q_{P_0}$、$\mu_{P_0}$ 和 $g_{P_0}$，并对它们作波动。有兴趣的读者可参阅该包的手册及示例说明文档。Sampford抽样方法（Sampford，1967）实现了 $P^{s\mathrm P}$。然而，当比值 $n/N$ 接近0或1时，在这里即 $n/N$ 不等于 $10^{-3}$ 时，这一接受—拒绝算法通常需要很长时间才能成功。为绕开这一问题，我们用Pareto抽样来近似 $P^{s\mathrm P}$，见 Bondesson 等（2006）第5节中的算法2。我们按照第29.2.2节描述的方法实现 $P^{s\Pi}$，并作少量调整，以处理这样的情况：对某个 $1\le v\le3$，$\sum_{i=1}^NK_{ii}I\{V_i=v\}$ 可能不是整数，或不能整除 $\sum_{i=1}^NI\{V_i=v\}$。

结果汇总于表29.1。我们重点考察TMLE的经验覆盖率、经验方差及方差估计的均值。除一个等于94\%的经验覆盖率外，其余都高于95\%。在每种情形中，方差估计的均值都大于对应的经验方差，表明我们实现了对 $\sigma_1^2$ 的保守估计。关于方差，在 $P_{0,1}$ 下，$P^{s\mathrm P}$ 与 $P^{s\Pi}$ 的表现相似，并且都略优于SRS。这与预期一致，因为 $P_{0,1}$ 下给定 $(A,W)$ 时 $Y$ 的条件标准差所形成的反差，正会产生这种结果。在 $P_{0,2}$ 下，$P^{s\mathrm P}$ 与 $P^{s\Pi}$ 的表现也相似，并且显著优于SRS。这同样符合预期，因为给定 $(A,W)$ 时 $Y$ 的条件标准差所形成的反差，在 $P_{0,2}$ 下比在 $P_{0,1}$ 下更强。最后，在 $P_{0,3}$ 下，$P^{s\mathrm P}$ 优于SRS，$P^{s\Pi}$ 又略优于 $P^{s\mathrm P}$。这仍与预期一致，因为给定 $(A,W)$ 时 $Y$ 的条件标准差形成了反差，与 $P_{0,2}$ 下相同，而且 $P_{0,3}$ 与 $P_{0,2}$ 下给定 $(A,W)$ 时 $Y$ 的条件均值不同。

% 表29.1：原书分章PDF第15页，9行×9数据列=81格。
% P0,1经验方差行的09、08等前导零是原印，全部保留。
\begin{table}[htbp]
\centering
\caption{模拟研究结果汇总}
\label{tab:ch29-results}
\begingroup
\small
\setlength{\tabcolsep}{3pt}
\renewcommand{\arraystretch}{1.13}
\begin{tabular}{@{}ll*{9}{c}@{}}
\toprule
&&\multicolumn{3}{c}{SRS}&\multicolumn{3}{c}{$P^{s\mathrm P}$，使用 $h_{n_0}^{\mathrm P}$}&\multicolumn{3}{c}{$P^{s\Pi}$，使用 $h_{n_0}^{\Pi}$}\\
\cmidrule(lr){3-5}\cmidrule(lr){6-8}\cmidrule(lr){9-11}
&&\multicolumn{3}{c}{$n$}&\multicolumn{3}{c}{$n$}&\multicolumn{3}{c}{$n$}\\
&&$10^3$&$5\times10^3$&$10^4$&$10^3$&$5\times10^3$&$10^4$&$10^3$&$5\times10^3$&$10^4$\\
\midrule
$P_{0,1}$&经验覆盖率&96.2\%&98.9\%&99.2\%&98.1\%&98.6\%&99.4\%&97.8\%&99.2\%&99.3\%\\
&经验方差&09&08&07&07&06&06&07&06&06\\
&估计方差&13&14&14&11&11&11&11&11&11\\
\midrule
$P_{0,2}$&经验覆盖率&94.0\%&98.9\%&99.2\%&98.9\%&99.9\%&99.1\%&98.4\%&99.4\%&99.3\%\\
&经验方差&129&104&102&44&41&44&49&42&42\\
&估计方差&171&196&200&85&86&87&86&85&86\\
\midrule
$P_{0,3}$&经验覆盖率&95.6\%&98.8\%&97.8\%&97.8\%&97.9\%&97.1\%&98.5\%&98.3\%&96.3\%\\
&经验方差&157&134&168&85&91&116&81&85&104\\
&估计方差&216&242&245&130&133&135&124&128&127\\
\bottomrule
\end{tabular}
\endgroup
\par\medskip
\begin{minipage}{\linewidth}
\footnotesize
上、中、下三组行分别对应 $P_{0,1}$、$P_{0,2}$ 和 $P_{0,3}$ 下的模拟。对每个子样本量 $n$ 及每种调查抽样设计，各组都报告置信区间的经验覆盖率 $B^{-1}\sum_{b=1}^BI\{\Psi^c(P_{0,j})\in I_{n,b}\}$、估计量经验方差的 $n$ 倍 $n[B^{-1}\sum_{b=1}^B\psi_{n,b}^{*2}-(B^{-1}\sum_{b=1}^B\psi_{n,b}^*)^2]$，以及估计量方差估计的 $n$ 倍的经验均值 $B^{-1}\sum_{b=1}^B\hat\sigma_{n,b}^2$。
\end{minipage}
\end{table}

